\documentclass[11pt,a4paper, titlepage=false]{scrreprt}

\usepackage[utf8]{inputenc}
\usepackage[T1]{fontenc}
\usepackage{lmodern}
\usepackage[english]{babel}
\usepackage{amsmath,amsfonts,amssymb}
\usepackage{cite}
\usepackage{hyperref}
\usepackage{graphicx}
\usepackage{url}
\usepackage{breakurl}
\usepackage{xcolor}
\usepackage{enumitem}
\usepackage[
  left=2.5cm,
  right=2.5cm,
  top=3cm,
  bottom=3cm,
]{geometry}

\hypersetup{
    pdfauthor  ={Muhammad Sawaiz Naveed},
    pdftitle   ={LLM-Driven Verilog Testbench Generation with Pyverilog-Based Early Error Localization},
    pdfsubject ={S6.ReKI.1 Research Project},
    pdfkeywords={LLM, Verilog, testbench generation, Pyverilog, error localization, hardware verification, LangGraph}
}

\usepackage{listings}
\DeclareOldFontCommand{\bf}{\normalfont\bfseries}{\mathbf}

%%%%%%%%%%%%%%%%%%%%%%%%%%%%%%%%%%%%%%%%%%%%%%%%%%%%%%%%%%%%%%%%%%%%%%%%%%%%%%%%%%%%%%%%%
%                                   Document
%%%%%%%%%%%%%%%%%%%%%%%%%%%%%%%%%%%%%%%%%%%%%%%%%%%%%%%%%%%%%%%%%%%%%%%%%%%%%%%%%%%%%%%%%

\begin{document}
\selectlanguage{english}
\title{LLM-Driven Verilog Testbench Generation\\with Pyverilog-Based Early Error Localization}
\author{Muhammad Sawaiz Naveed}
\date{\today}
\maketitle
\vspace{1cm}

\begin{tabular}{r l}
Type of work: & Research Project \\
Student:      & Muhammad Sawaiz Naveed \\
Supervisor:   & Bing Wen \\
Topic ID:     & S6.ReKI.1 \\
Department:   & Technische Universität Ilmenau \\
Semester:     & Summer Semester 2025/2026 \\
Deadline:     & September 1, 2026 \\
\end{tabular}

\vspace{3mm}

\begin{tabular}{l p{0.75\textwidth}}
\\
\textbf{Keywords:} & Large Language Models; Verilog; Testbench Generation; Pyverilog; Error Localization; Hardware Verification; LangGraph \\
\end{tabular}

\vspace{8mm}

%%%%%%%%%%%%%%%%%%%%%%%%%%%%%%%%%%%%%%%%%%%%%%%%%%%%%%%%%%%%%%%%%%%%%%%%%%%%%%%%%%%%%%%%%

\section*{Motivation}

The increasing complexity of digital hardware systems has made verification one of the most
resource-intensive phases of the design process. According to industry studies, verification
tasks alone can consume up to 60\% of the total engineering effort in modern IC
projects~\cite{wilson2022}, and a substantial share of that effort goes into writing and
maintaining testbenches. Large Language Models (LLMs) have recently demonstrated promising
capabilities in generating Hardware Description Language (HDL) artefacts --- both design
RTL~\cite{liu2023verilogeval,thakur2023autochip} and, more recently, complete
testbenches~\cite{qiu2024autobench} --- directly from natural-language descriptions. This
raises the prospect of significantly accelerating the hardware verification flow through
automation.

However, LLM-generated Verilog is fundamentally unreliable. Verilog demands strict adherence to
parallel semantics, clock-domain discipline, synthesizability constraints, and precise signal
timing. LLMs, trained predominantly on sequential software code, frequently produce code that
is \emph{syntactically valid} yet \emph{functionally wrong}: the simulator compiles the code
without errors while the underlying behaviour is incorrect. In the testbench setting, common
failure patterns include wrong port bindings between the testbench driver and the
Device-Under-Test (DUT), incorrect or missing \texttt{\$fdisplay} statements that prevent the
checker from observing outputs, mismatched bit widths between stimulus and DUT ports, wrong
sensitivity lists in clocked drivers, and incomplete or off-by-one scenario coverage.

The conventional remedy is to run the LLM output and inspect the simulation result. While
effective, this is a slow feedback loop: the testbench must compile, the DUT must be
instantiated correctly, and the run must complete before any feedback is available. Recent work
such as AutoBench~\cite{qiu2024autobench} shows that a multi-stage LLM pipeline can substantially
improve testbench pass rates, but it still relies almost exclusively on the compiler and the
final simulation output to detect problems. What is missing is a \emph{lightweight, pre-simulation
analysis layer} that can inspect the generated testbench (and its interaction with the DUT),
localise common functional errors immediately, and feed structured, signal-level guidance back
to the LLM --- without depending on a full simulation run inside every repair iteration.

This project addresses that gap. By embedding Pyverilog's static-analysis capabilities (AST
parsing, dataflow tracing, port-binding and sensitivity-list checks) inside a graph-based LLM
workflow, we aim to generate Verilog testbenches that are more robust by construction, and to
localise and repair the residual errors \emph{early}.

%%%%%%%%%%%%%%%%%%%%%%%%%%%%%%%%%%%%%%%%%%%%%%%%%%%%%%%%%%%%%%%%%%%%%%%%%%%%%%%%%%%%%%%%%

\section*{Task Description}

The goal of this project is to design, implement, and evaluate a \emph{graph-based LLM workflow
for Verilog testbench generation} that integrates Pyverilog-based early error localisation as
a first-class node in the pipeline. The pipeline takes a natural-language circuit description
and a golden DUT as input and produces a working testbench, together with a structured trace of
every node executed and every error localised and repaired along the way.

The focus of the work is on the \emph{pipeline architecture and its error-handling behaviour},
rather than on achieving the highest possible benchmark score. The objective is a clean,
reproducible, graph-based design in which each stage --- generation, static analysis, error
reasoning, repair, and evaluation --- is an explicit node with logged inputs and outputs.

\medskip
\noindent The project pursues the following research questions:

\begin{enumerate}[label=\textbf{RQ\arabic*.}, leftmargin=*, itemsep=4pt]

    \item \textit{What categories of functional errors appear most frequently in LLM-generated
    Verilog testbenches, and which of these are detectable without full simulation?}

    \item \textit{To what extent can Pyverilog's AST and dataflow analysis (port bindings,
    sensitivity lists, dataflow consistency between testbench and DUT) narrow down testbench
    errors prior to simulation?}

    \item \textit{Can an LLM, informed by Pyverilog analysis results, effectively localise and
    repair testbench errors, and how does this compare to using only compiler/simulator feedback
    as in prior work?}

    \item \textit{What is the cost--quality tradeoff of combining lightweight static analysis
    with LLM reasoning versus relying solely on compiler and simulation feedback for repair?}

\end{enumerate}

\medskip
\noindent The concrete outcomes of the project are:

\begin{itemize}
    \item A \textbf{LangGraph-based testbench generation pipeline} with explicit nodes for
    circuit classification, specification extraction, scenario generation, driver/checker
    generation, Pyverilog static analysis, LLM error reasoning, deterministic standardisation,
    and Icarus Verilog evaluation.

    \item A \textbf{reusable Pyverilog analysis module} that converts AST and dataflow outputs
    into structured, LLM-readable error summaries focused on testbench-DUT interaction
    (port-binding mismatches, undriven signals, sensitivity-list errors, missing observation
    statements).

    \item A \textbf{deterministic Python standardiser} for sequential testbenches that inserts
    missing \texttt{\$fdisplay} statements and normalises clocking, replacing the fragile
    LLM-based standardisation step used in prior work~\cite{qiu2024autobench}.

    \item A \textbf{categorised error taxonomy} of testbench errors observed during the study,
    annotated with frequency and Pyverilog detectability.

    \item An \textbf{empirical comparison} of three feedback strategies for the repair loop ---
    compiler/simulator-only, Pyverilog-only, and a hybrid Pyverilog~+~LLM reasoner --- across
    combinational and sequential designs, with per-stage failure attribution made possible by
    the LangGraph logging layer.
\end{itemize}

%%%%%%%%%%%%%%%%%%%%%%%%%%%%%%%%%%%%%%%%%%%%%%%%%%%%%%%%%%%%%%%%%%%%%%%%%%%%%%%%%%%%%%%%%

\section*{Scientific Approach}

\subsection*{Sources and Related Work}

The project builds on the following bodies of work:

\begin{itemize}
    \item \textbf{LLM-based Verilog generation:} VerilogEval~\cite{liu2023verilogeval} provides
    the standard benchmark (156 Verilog problems from HDLBits); AutoChip~\cite{thakur2023autochip}
    and ChipGPT~\cite{chang2023chipgpt} demonstrate LLM-driven RTL generation pipelines.

    \item \textbf{Testbench and verification automation:} AutoBench~\cite{qiu2024autobench}
    introduces a multi-stage LLM pipeline for testbench generation, achieving a 57\% improvement
    in pass rate over direct LLM prompting. Its self-enhancement loop (scenario checking,
    auto-debug, reboot) informs our repair loop design.

    \item \textbf{Pyverilog:}~\cite{takamaeda2015pyverilog} provides the AST parser, dataflow
    analyser, and control-flow analyser used for static analysis of Verilog code.

    \item \textbf{LLMs for formal and functional verification:}~\cite{orenes2023llmformal}
    and~\cite{zhang2023llm4dv} explore LLM-assisted formal verification and hardware test
    stimulus generation, providing context for the broader verification landscape.
\end{itemize}

The primary dataset is the public VerilogEval benchmark~\cite{liu2023verilogeval} (156 problems
from HDLBits), which provides natural-language descriptions, golden RTL, and reference
testbenches in a uniform format. A development split (20\%) is used for prompt tuning and
error-taxonomy bootstrapping; the held-out split (80\%) is reserved for final evaluation. If
the supervisor provides an additional dataset, it is added on top of this baseline.

\subsection*{Methods and Implementation}

The pipeline is structured as a LangGraph state machine with six explicit stages. The accent of
the work is on the \emph{shape} of the graph (nodes, edges, conditional transitions, repair
loop) rather than on raw benchmark numbers.

\begin{enumerate}[itemsep=3pt]
    \item \textbf{Multi-stage Testbench Generation.} The LLM produces the testbench in small,
    auditable steps rather than one shot: (a) classify the DUT as combinational or sequential;
    (b) extract a structured JSON specification of ports, behaviour, and timing; (c) generate
    a named scenario list; (d) generate Verilog driver code that exercises those scenarios; and
    (e) generate a Python checker. This staged structure follows the recipe shown to work in
    AutoBench~\cite{qiu2024autobench} and gives the static-analysis layer a well-formed object
    to inspect.

    \item \textbf{Pyverilog Static Analysis.} The generated testbench and the golden DUT are
    parsed together. Pyverilog is used to (i) verify port bindings between the testbench
    instantiation and the DUT signature, (ii) detect undriven inputs and unobserved outputs,
    (iii) check sensitivity lists in sequential drivers, and (iv) confirm that
    \texttt{\$fdisplay} (or equivalent) statements are present for every output. No simulator is
    invoked at this stage.

    \item \textbf{LLM Error Reasoning.} The structured Pyverilog report and the original
    specification are passed to an LLM, which produces a localised error report
    (error type, affected signal, suggested fix). Where it can, the reasoner cites the line or
    block in the testbench that is responsible.

    \item \textbf{Deterministic Standardiser.} A small Python pass walks the testbench AST and
    inserts missing \texttt{\$fdisplay} statements and normalises clocking conventions. This
    replaces the LLM-based standardisation step in prior work, which is fragile on a mechanical
    transformation~\cite{qiu2024autobench}.

    \item \textbf{Repair Loop.} If errors remain, the error report is fed back to the
    generation node and the testbench is regenerated. The loop runs up to $N_{\max} = 3$
    iterations; an oscillation detector breaks early if the same error report repeats. The full
    LangGraph trace (which node produced which artefact at which iteration) is logged.

    \item \textbf{Evaluation with Icarus Verilog.} The final testbench is compiled and run with
    Icarus Verilog~\cite{blocklove2023chipchat} against the golden DUT, with three checks:
    Eval0 (compiles), Eval1 (passes against golden RTL), and Eval2 (catches LLM-generated
    mutants of the DUT, following AutoBench's mutant-pass protocol).
\end{enumerate}

\noindent The LangGraph state machine makes every transition explicit and logged, supporting
reproducibility and per-node failure attribution --- which itself becomes part of the
contribution.

\medskip
\noindent\textbf{Tooling.} Pyverilog~\cite{takamaeda2015pyverilog} is the primary static
analyser, with a Verible fallback for files Pyverilog cannot parse. The LLM tier is the free
Claude API (Sonnet for code generation and reasoning, Haiku for cheap classification nodes), in
line with the project's emphasis on pipeline structure rather than maximising accuracy. All
code-generation calls use temperature~0 for reproducibility.

\subsection*{Evaluation Metrics}

The primary goal of the evaluation is to characterise how each node of the pipeline contributes
to overall behaviour, not to chase peak benchmark numbers.

\begin{itemize}
    \item \textbf{Eval0:} Testbench compilation pass rate (Icarus Verilog).
    \item \textbf{Eval1:} Testbench passes against the golden DUT.
    \item \textbf{Eval2:} Testbench correctly distinguishes the golden DUT from LLM-generated
    mutants (AutoBench's mutant-pass protocol).
    \item \textbf{Error precision / recall:} Fraction of Pyverilog-flagged errors that are real,
    and fraction of real errors Pyverilog catches, measured against a hand-annotated dev-set
    subset.
    \item \textbf{Per-node failure attribution:} Distribution of failures over the six pipeline
    stages, exposing where the pipeline is fragile.
    \item \textbf{Iterations to pass:} Distribution of repair-loop iterations required.
    \item \textbf{Cost per module:} LLM tokens and wall-clock time per generated testbench.
\end{itemize}

\noindent Results are reported separately for combinational (CMB) and sequential (SEQ)
circuits, as the latter are substantially harder for LLMs~\cite{qiu2024autobench}. The
ablation conditions exercise the value of each feedback strategy: (a) baseline (one-shot
generation, no repair), (b) compiler/simulator-feedback-only repair, (c) Pyverilog-only repair,
(d) LLM-reasoning-only repair, and (e) the full hybrid pipeline.

%%%%%%%%%%%%%%%%%%%%%%%%%%%%%%%%%%%%%%%%%%%%%%%%%%%%%%%%%%%%%%%%%%%%%%%%%%%%%%%%%%%%%%%%%

\section*{Time Schedule}

\vspace{0.7em}
\begin{tabular}{|c | p{0.75\textwidth}|}
\hline
\textbf{Time Period} & \textbf{Tasks} \\
\hline
Weeks 1--2 (May 2026) &
    Literature review (AutoBench, VerilogEval, Pyverilog); dev environment setup
    (Python, Claude API free tier, Icarus Verilog, LangGraph); Pyverilog/Verible
    smoke test on 20 sample modules; report skeleton committed \\
\hline
Weeks 3--6 (May--Jun 2026) &
    LangGraph skeleton and state schema; testbench-generation stages 1a--1c
    (classification, spec extraction, scenario list) on combinational circuits;
    Icarus Verilog wiring for Eval0/Eval1 \\
\hline
Weeks 5--9 (Jun 2026) &
    Pyverilog static-analysis module (port-binding, sensitivity, dataflow);
    LLM error-reasoning node; error-taxonomy bootstrap on a 30--50-module dev
    subset; first ablation numbers \\
\hline
Weeks 10--13 (Jun--Jul 2026) &
    Repair loop with oscillation detection; deterministic \texttt{\$fdisplay}
    standardiser; sequential-circuit support; Eval2 mutant generation; full
    pipeline integration \\
\hline
Weeks 14--16 (Jul--Aug 2026) &
    Ablation experiments on the held-out test set; per-node failure attribution;
    cost/latency analysis; complete first draft of the final report \\
\hline
Weeks 17--20 (Aug--Sep 2026) &
    Supervisor feedback round; revision and proofreading; artefact preparation
    (code, prompts, reproducibility scripts) \\
\hline
September 1, 2026 & \textbf{Submission of final report (PDF upload).} \\
\hline
\end{tabular}

%%%%%%%%%%%%%%%%%%%%%%%%%%%%%%%%%%%%%%%%%%%%%%%%%%%%%%%%%%%%%%%%%%%%%%%%%%%%%%%%%%%%%%%%%

\vspace{5mm}
\section*{References}
{\small
\begingroup
\let\chapter\section
\bibliographystyle{ieeetr}

\begin{thebibliography}{9}

\bibitem{wilson2022}
H.~Foster,
``The 2022 Wilson Research Group Functional Verification Study,''
Siemens, Tech. Rep., Dec. 2022.
\url{https://blogs.sw.siemens.com/verificationhorizons/2022/12/12/part-8-the-2022-wilson-research-group-functional-verification-study/}

\bibitem{liu2023verilogeval}
M.~Liu, N.~Pinckney, B.~Khailany, and H.~Ren,
``VerilogEval: Evaluating Large Language Models for Verilog Code Generation,''
in \textit{Proc. IEEE/ACM Int. Conf. Computer-Aided Design (ICCAD)}, 2023.

\bibitem{qiu2024autobench}
R.~Qiu, G.~L.~Zhang, R.~Drechsler, U.~Schlichtmann, and B.~Li,
``AutoBench: Automatic Testbench Generation and Evaluation Using LLMs for HDL Design,''
in \textit{Proc. ACM/IEEE Int. Symp. Machine Learning for CAD (MLCAD)}, 2024.
arXiv:2407.03891.

\bibitem{takamaeda2015pyverilog}
S.~Takamaeda-Yamazaki,
``Pyverilog: A Python-Based Hardware Design Processing Toolkit for Verilog HDL,''
in \textit{Applied Reconfigurable Computing (ARC)}, 2015.

\bibitem{thakur2023autochip}
S.~Thakur, J.~Blocklove, H.~Pearce, B.~Tan, S.~Garg, and R.~Karri,
``AutoChip: Automating HDL Generation Using LLM Feedback,''
arXiv:2311.04887, 2023.

\bibitem{blocklove2023chipchat}
J.~Blocklove, S.~Garg, R.~Karri, and H.~Pearce,
``Chip-Chat: Challenges and Opportunities in Conversational Hardware Design,''
in \textit{Proc. ACM/IEEE Int. Symp. Machine Learning for CAD (MLCAD)}, 2023.

\bibitem{chang2023chipgpt}
K.~Chang, Y.~Wang, H.~Ren, M.~Wang, S.~Liang, Y.~Han, H.~Li, and X.~Li,
``ChipGPT: How Far Are We From Natural Language Hardware Design,''
arXiv:2305.14019, 2023.

\bibitem{orenes2023llmformal}
M.~Orenes-Vera, M.~Martonosi, and D.~Wentzlaff,
``Using LLMs to Facilitate Formal Verification of RTL,''
arXiv:2309.09437, 2023.

\bibitem{zhang2023llm4dv}
Z.~Zhang, G.~Chadwick, H.~McNally, Y.~Zhao, and R.~Mullins,
``LLM4DV: Using Large Language Models for Hardware Test Stimuli Generation,''
arXiv:2310.04535, 2023.

\end{thebibliography}

\endgroup
}
\end{document}
